\scanpage{15}
Schémas (finis) étales sur un corps $k$ : la catégorie est équivalente : la catégorie des $\pi$-ensembles (finis), $\pi=\operatorname{Gal}(\bar k/k)$.

\uline{Application} Soit $X$ un schém. sur $k$, Supposons d'abord que $X_{\bar k}$ est loc. connexe [p. ex. loc. noeth, ou ens des composantes irréductibles loc. fini ; \unsure{tout est ok|tout est OK} si $X$ est loc. de t.f. sur $k$]. Alors $I=\pi_0(X_{\bar k})$ est un ens. discret, et on a un \deleted{\unsure{$\varphi$|f}} morphisme de $\bar k$-schémas

\diagram{\begin{tikzcd}[ampersand replacement=\&]
X_{\bar k}\arrow[d]\\
I_{\bar k}
\end{tikzcd}}
(il est universel pour les morphismes de $X$ dans des schémas constants). D'ailleurs $\pi$ opère sur $\deleted{I_{\bar k}} I$, donc $I_{\bar k}$ provient d'un \deleted{$I$} \unsure{sur|\deleted{sur}} schém. étale $T$ sur $k$ [si on prouve que $\pi$ opère continûment -- ce sera le cas en tous cas si $X$ est quasi-compact].

\margin{i.e \deleted{tte} composante connexe de $X_{\bar k}$ est \unsure{déjà|déja} définie sur une \inserted{\unsure{ss-|s.-|sép.}}ext. finie de $k$. \uline{C'est vrai} \uline{si} $X$ \uline{qu. cpct}, \deleted{\unsure{mais|aussi}} \deleted{\illegible{2}}}

% slanted
Alors la théorie de la descente nous prouve que $X\simeq X_{\bar k}/\pi$, $T=I_{\bar k}/\pi$, (dans catégorie des schémas) donc on trouve

\diagram{\begin{tikzcd}[ampersand replacement=\&]
X\arrow[d,"f"]\\
T
\end{tikzcd}}
$T$ étant étale sur $k$, \deleted{et} On \unsure{a|â|à} que $f$ est surjectif et à fibres géom connexes -- et ceci suffit à caractériser $X$. On peut aussi décrire comme le morphisme universel de $X$ dans un schém. étale sur $k$.

\margin{$T$ est appelé le schém. des composantes connexes géom. de $X$. Un pt à val. dans un $S'/k$ ``est'' une partie ouverte et fermée de $X_{\unsure{\xi|S'|S}}$, \unsure{/|,} qui \unsure{est|\deleted{est}}}

\diagram{\begin{tikzcd}[ampersand replacement=\&,column sep=large]
X\arrow[d] \& X'\arrow[d,dash]\\
T\arrow[d] \& T'\arrow[d,dash]\\
S=\operatorname{Sp} k\arrow[r,dash] \& S'
\end{tikzcd}}
\uline{Généralisation} On suppose $X$ quasi-compact, ou somme de schémas quasi cpcts : On trouve un $T$ \inserted{affine} pro étale, ayant les \unsure{\={m}s|\~{m}s} propriétés$\scriptscriptstyle\wedge$ qu'avant. Le foncteur qu'il représente est un peu plus subtil à décrire....

% archive sheet number: 14
